% Scientific Reports Article Discussion: continuous prose without
% subheadings, no display items; the final paragraph states what is
% established, what is a design argument and what is open. Derived from
% 04_discussion_CM.tex. Relocated rather than deleted:
%   - PID literature positioning and the differentiation from Liang et al.
%     -> Methods, "Relation to prior information-decomposition work"
%   - the noise-aware estimation argument -> Methods (already present)
%   - the fidelity spectrum of cross-domain tasks -> Supplementary Fig.
%     fig:fidelity_spectrum legend and Supplementary Note 3
%   - glioma/ECG instantiation detail and the CPTAC controls -> Methods and
%     Supplementary Note 1
%   - the extended limitations block -> Supplementary Note 4

\section*{Discussion}

Computational pathology has converged on a single output for mutation
prediction, a probability per gene, and on two assumptions: that it is
comparable across genes, and that improvement on the
cohort mean is the figure of merit. Our results challenge both. One
cohort, encoder and backbone yield qualitatively different gene-level
behaviour: an encoder-ceiling null, a label-space mismatch and an
estimation floor are three failures a leaderboard reports as one
mediocre number, calling for three different remedies. The
Conclusive/Inconclusive flag is therefore not cosmetic: the cohort
metric does not translate uniformly into per-slide actionability, so
downstream clinical use must be gene-conditional.

Reading the results this way also changes what counts as the
contribution. The pattern-informed channel does not win the cohort-mean
benchmark, and with the fusion signs disagreeing across seeds neither
magnitude nor direction can be defended (\cref{tab:luad_contrasts}). The
one contrast that survives says what the vocabulary omits, not what
fusion adds. An instrument that reports this, rather than the largest
cell of the table, is the result; the alternative reading, in which
the $0.03$ gain of \fcmil{} over \piabmil{} on \textit{STK11} is a
finding, is exactly the kind of claim the audited literature makes at
the same effect size and that the bootstrap here shows to be
unsupported.

What pattern-informed aggregation wins is not accuracy but the argument
with the pathologist. A model that names a positive
lepidic\,$\times$\,solid interaction invites a textbook reply about the
invasive-mucinous proxy, and the conditional test then settles whether
the reply holds (\cref{tab:kras_coactivation}). A model that returns
KEAP1 at chance on both channels is transparent about what the slide
cannot answer, and points the next investigation toward NRF2
immunostaining, not another loss function.\cite{rudin2019stop}

What lifts this from an empirical grouping to a portable claim is where
the categories come from: a taxonomy that sorts genes by \auroc{}
describes one cohort and travels no further. These categories are
instead the four ways two sources can distribute information about a
target---an exhaustive partition, not an observed one---and so
the taxonomy makes checkable commitments: the
non-additive aggregator must win on exactly the genes with a recoverable
synergistic atom, lose elsewhere, and---binding
hardest---win \emph{where the named pattern pair co-activates}, not
on average. That last commitment is the one the cohort could test, and
on KRAS it returned empty; a taxonomy whose predictions cannot fail
would not be worth porting. The objection this must answer is that the regimes restate
what is already known: that KEAP1 acts through a biochemical axis, that
STK11 acts through the immune microenvironment. Recovering those is
\emph{calibration}, and an instrument earns trust
by returning the known answer where the answer is known. The predictive
content lies where prior knowledge does not reach: the additive sub-form
(ii-a) is an empty cell of \cref{eq:pid} in LUAD, nameable because the
lattice enumerates it and not because a case was observed; its
completeness rests on the algebra, not on any cohort. What the framework
predicts is not the mutation but which remedy will move it.

That changes how the field decides what to build next. Under a cohort
mean, ``what should we add?'' has one answer for every gene, usually
more parameters or more slides. Under \cref{eq:pid} it is
atom-conditional: a redundant atom binds on the encoder, a synergistic
one on the aggregator's hypothesis space, an empty vocabulary-side atom
on the \emph{source vocabulary}, and a non-empty but unestimable one on
the \emph{cohort}.

The sharper consequence is that untranslatability becomes a
first-class, actionable result---provided it is asserted at the level
where it was established, the qualification the framework most
tempts one to skip. The
verdict thresholds a mean \auroc{}, a property of the cohort and of no
slide, so reading a low mean as ``this mutation cannot be seen in this
tissue'' is a level error, and a consequential one, because it discards
exactly the cases where the morphology does speak.
Untranslatability proper is the narrower claim
that the source-side atoms are empty at the level being discussed. When a
gene meets it, the principled output is not a low probability presented
as a confident call, but the statement that the mutation cannot be said
in morphological words, and---from the single-channel conditions---why.
The operational rule
is restrictive: an Inconclusive gene licenses no
diagnostic use, not a hedged call and not a ranked shortlist, until the
deficit the decomposition names is repaired.

Two facts about the evidence are structural rather than incidental, and
we state them as such. The first is that regime~(ii) is occupied by no
target in any cohort here. The taxonomy is a partition of
\cref{eq:pid}, so the cell exists whether or not a cohort has filled
it, and the rule states what would: \piabmil{}$-$B2 excluding zero with
\fcmil{}$\approx$\piabmil{} for the additive sub-form, or an
\fcmil{}$-$\piabmil{} contrast concentrating in the co-activation
stratum for the non-additive one. Until a cohort supplies that
observation, (ii) is the instrument's one untested prediction, and the
accurate description of the taxonomy is two regimes observed and one
specified. A $63$-subject glioma pilot with expert-scored VASARI points
towards (ii-a) for 1p/19q---concat fusion $+0.14$ over the encoder
(\cref{tab:glioma_domainb})---but its fold spread ($\pm0.21$) and a sex
control returning a fusion gain of comparable size make it an
illustration of where the cell would be filled, not an observation that
fills it. The second is scarcity, and it is the condition the
instrument is built for rather than a limitation it suffers. All
primary numbers are cross-validation on a single retrospective cohort
that the one external cohort we score does not reproduce off-the-shelf;
at $458$ slides from $444$ patients TCGA-LUAD is one to two orders of
magnitude smaller than the proprietary cohorts of the foundation-scale
line, and every public histology-to-genotype cohort is of that order.
At that size the design resolves $0.04$--$0.11$ \auroc{} per target, and
detecting a $0.03$ aggregator gain, mid-range for this literature,
would need roughly $9{,}000$ patients patient-preserved and $13{,}400$
site-preserved (\cref{tab:design_power}). No single institution holds
that, which makes interaction-level mechanistic claims
multi-institutional by construction; what a single-institution cohort
can honestly report is the number itself, and that is what the
instrument returns where a cohort mean would return a false assignment.
The encoder bounds the numbers too, and we measured that rather than
assuming it: a permutation null returns chance and tumour-versus-normal
$0.98$ while \textit{TP53} stops at $0.67$, so the pipeline is
calibrated but its ceiling is low, and a stronger encoder could move
genes between regimes. Structural portability and the stability of a
given gene's assignment are therefore different claims: because the
atoms and the verdict are population expectations, an assignment is a
property of the triple (vocabulary, encoder, cohort) rather than of the
gene, and calling one near-intrinsic would need cross-cohort concordance
a single-cohort study cannot supply.

Three things are established here, one is a design argument, and two
are open. Established: that the four-condition protocol with its
noise-aware rule returns a reading on every target it is given, across
three vocabularies and three encoders with nothing in common---certifying
equivalence where a cohort is large enough (the electrocardiogram) and,
under the scarcity of every public cohort, returning a warranted
null with the sample size that would resolve it, in place of an
assignment; that the synergy signature is falsifiable and survived its
one available test by returning empty; and that a predictor's collapse
under cohort transfer can be localised to the transported decision
function rather than to the tissue. The design argument is that the
vocabulary the verdict is stated in is what would let a pathologist
contest it; the per-slide tooling released with the code (Methods) is
not evaluated here and no claim about reader behaviour is made. Open
are whether any per-target regime assignment reproduces across cohorts,
which no single-cohort study can settle, and whether any target occupies
regime~(ii), which no cohort here does. The instrument thus separates
the two things a low cohort mean conflates: a target whose deciding
contrast is wider than the margin, which more patients would resolve and
for which it returns the number (Methods, \cref{eq:design_n,eq:equivalence_n}),
and a target whose every condition sits at chance, which no cohort size
would change. A cohort-size recommendation is attached to a contrast,
never to a target's \auroc{}; that is what reading regimes instead of
scores buys. The actionable conclusion is therefore not that this
predictor outperforms the published competition---it does not, and is
not expected to at this cohort size---but that cross-domain medical
prediction should be reported with per-target Conclusive/Inconclusive
flags, and that a taxonomy grounded in a partial information
decomposition is a more rigorous scaffold for cross-paper comparison
than the cohort average now in use. Being a decomposition rather than a
naming convention, it applies unchanged to any task pairing an expert
ontology with a learned encoder.
